% Proposed A.86. Source only until independently reviewed and privately built.
% Combines the three frozen checkpoint27 cusp notes without changing them.
\statement{A.86}{Proved}{A canonical full cusp word compares the dihedral and relative Alexander units}
Retain the actual negative-\(\widehat{GT}_0\) data of A.85,
and put \(z=(xy)^{-1}\), \(G_v=F\rtimes\langle v\rangle\).
Let \(d=d_\alpha\) be the actual full unit from A.45:
in the universal dihedral quotient
\[
 \pi(x)=a,\qquad\pi(y)=aT,\qquad\pi(u)=T^k,
 \qquad d=2k+1,\qquad k\in2\widehat{\mathbb Z}.
 \tag{A.86.1}
\]
In particular \(d_2\equiv1\pmod4\). There is a unique
\(h_0\in F'\) satisfying
\[
             \tau(z)=h_0^{-1}z^{-1}h_0,
\]
and a unique \(q\in F\) such that
\[
 h_0\tau=q\tau q^{-1},\qquad q_{\mathrm{ab}}=(t,-t)
       \quad\text{for some }t\in\widehat{\mathbb Z}.
 \tag{A.86.2}
\]
Define actual full-group elements
\[
 \begin{gathered}
 e=(xu)v,\qquad b=q^{-1}eq,\qquad
 \Omega_\alpha=(x^dv)b^{-d}\in F,\\
 \sigma=\operatorname{Inn}(h_0)\tau,\qquad
 g=e(qvq^{-1})=xu\,v(q)q^{-1}\in F.
 \end{gathered}
 \tag{A.86.3}
\]
Then
\[
 \begin{gathered}
 e^2=xy,\qquad\tau(b)=b^{-1},\qquad
 \sigma(g)g^{-1}=z,\\
 a_{\mathrm{reg}}:=1-4t\in\widehat{\mathbb Z}^{\times},\qquad
 (\Omega_\alpha)_{\mathrm{ab}}
   =\left(\frac{d-a_{\mathrm{reg}}}{2},
                  -\frac{d-a_{\mathrm{reg}}}{2}\right).
 \end{gathered}
 \tag{A.86.4}
\]
The element g is exactly C.44's unique solution for the normalized
regular involution \(\sigma\); \(a_{\mathrm{reg}}\) is its
image in \(F/\langle\!\langle z\rangle\!\rangle\), with
generator the image of x. The full criterion is
\[
 \begin{aligned}
 \Omega_\alpha=1
 &\quad\Longleftrightarrow\quad
 x^du^{-1}y^du\sim_F(xy)^d\\
 &\quad\Longleftrightarrow\quad
 \alpha\text{ is conjugate to }c\text{ in }\widehat{GT}_0.
 \end{aligned}
 \tag{A.86.5}
\]
If \(\alpha\in\widehat{GT}\), the last condition is also
equivalent to conjugacy in \(\widehat{GT}\). These last group
membership assertions use exactly A.39, with A.6 and A.36 for the
respective descents, and A.45 for the fixed multiplier d.

When (A.86.5) holds, the original cusp equation has the specified
conjugator
\[
       x^du^{-1}y^du=q^{-1}(xy)^dq.
 \tag{A.86.6}
\]
It follows that the canonical Fox element of A.85 is a unit at every
coefficient prime. No converse from that Fox-unit condition alone is
asserted here. Neither \(a_{\mathrm{reg}}=d\) nor
\(\Omega_\alpha=1\) is proved by this statement.

\begin{proof}
\emph{The actual half rotation and its full centralizer.}\par
The identities \(v(x)=u^{-1}yu\) and \(v(u)=u^{-1}\) give
\[
 e^2=xu\,v(xu)=xy=z^{-1},\qquad
 v(z)=(xu)^{-1}z(xu).
 \tag{A.86.7}
\]
The actual isomorphism \(G_v\to F\rtimes\langle\iota\rangle\)
which fixes F and sends v to \(u^{-1}\iota\) sends e to
\(x\iota\). Its target is the full free profinite product
\(C_2(\iota)\amalg\widehat{\mathbb Z}(x)\).
Replacing x by \(x\iota\) is an invertible generator change,
so it also has the free factors \(C_2(\iota)\) and
\(\overline{\langle e\rangle}\). This is a description of
the actual ambient group; it is not a generation assertion for x and v.

The primitive cusp xy has full centralizer
\(C_F(xy)=\overline{\langle xy\rangle}\), as in C.5;
indeed \((xy,y)\) is a full free basis. The odd element e centralizes
it. An odd centralizing element differs from e by an even centralizing
element. Therefore
\[
 C_{G_v}(xy)=
 \overline{\langle xy\rangle}\sqcup
 \overline{\langle xy\rangle}e
 =\overline{\langle e\rangle}\simeq\widehat{\mathbb Z}.
 \tag{A.86.8}
\]
Injectivity of the displayed procyclic parametrization follows from
\(e^2=xy\): if \(e^a=1\), then \((xy)^a=1\), whence
\(a=0\). The same centralizer formula holds with xy replaced by
\((xy)^\lambda\) for any full unit \(\lambda\), since it
generates the same procyclic subgroup.

Put \(d_\lambda=x^\lambda v\). Then
\[
 d_\lambda^2=x^\lambda u^{-1}y^\lambda u,
 \qquad\tau(d_\lambda)=d_\lambda^{-1}.
 \tag{A.86.9}
\]
For every \emph{specified} \(s\in F\) one has
\[
 d_\lambda^2=s e^{2\lambda}s^{-1}
       \quad\Longleftrightarrow\quad
 d_\lambda=s e^\lambda s^{-1}.
 \tag{A.86.10}
\]
If the squares agree, \(s^{-1}d_\lambda s\) centralizes its
square and hence belongs to the group in (A.86.8). Write it as
\(e^a\). The equality \(2a=2\lambda\) in
\(\widehat{\mathbb Z}\) gives \(a=\lambda\). Thus the
conjugator is unchanged, rather than replaced by an unrelated
square-root conjugator. Conjugacy in \(G_v\) and in F are also
equivalent for these two elements: multiply an odd conjugator on
the right by \(e^{-1}\), which centralizes \(e^\lambda\)
and changes its parity. Under the actual free-product identification,
the full cusp condition is consequently
\[
       x^\lambda u^{-1}\iota\sim (x\iota)^\lambda.
 \tag{A.86.11}
\]
All of these are statements in the full groups.

\emph{The uniquely normalized actual cusp reflection.}\par
The regular pants property of \(\tau\) supplies an actual
\(h\in F\) with \(\tau(z)=h^{-1}z^{-1}h\).
The odd lift \(r=h\tau\) reverses z. Its square belongs to
\(C_F(z)=\overline{\langle z\rangle}\); it is both fixed
and inverted by r. Thus \(r^2=(r^2)^{-1}\), and torsion-freeness
of \(\overline{\langle z\rangle}\) gives \(r^2=1\). On
\(F_{\mathrm{ab}}=\widehat{\mathbb Z}^2\), \(\tau\)
acts by negative interchange. The equation \(r^2=1\) implies
\(h_{\mathrm{ab}}=(\ell,\ell)\) for some \(\ell\).
Thus
\[
                         h_0=z^\ell h\in F'
 \tag{A.86.12}
\]
has the same transport property. Any other transporter differs on
the left by \(z^a\), because the full centralizer of z is its
primitive procyclic group. Its abelianization changes by
\((-a,-a)\). Consequently exactly one transporter lies in F'.

The actual one-section theorem C.43 and \(C_F(\tau)=1\)
give the unique \(q\in F\) in (A.86.2). Abelianizing
\(h_0=q\tau(q)^{-1}\) gives
\(q_{\mathrm{ab}}=(t,-t)\). This construction uses the
existing actual cusp reflection; it does not assume a positive
peripheral conjugator, or membership of q in the marked subgroup H.

Now work in the actual full extension \(F\rtimes V\).
The element \(r=h_0\tau\) normalizes \(G_v\).
The square of \(rer^{-1}\) is \(e^{-2}\), so (A.86.8)
and cancellation of two in its exponent group give
\[
       rer^{-1}=e^{-1},\qquad
       \tau(q^{-1}eq)=(q^{-1}eq)^{-1}.
 \tag{A.86.13}
\]
This proves the asserted reversal of b. Its centralizer is the
conjugate torsion-free procyclic group \(\overline{\langle b\rangle}\).

\emph{The two normalized tests, with the same original conjugator.}\par
Suppose \(d_\lambda=s b^\lambda s^{-1}\), with \(s\in F\).
Applying \(\tau\) and using the two literal reversals gives
\[
       s^{-1}\tau(s)\in C_{G_v}(b^\lambda)
                            =\overline{\langle b\rangle}.
\]
Write it as \(b^a\). Since the left side is even and b has
v-parity one, \(a\in2\widehat{\mathbb Z}\). Then
\[
                   \tau(sb^{a/2})=sb^{a/2}.
\]
An element of \(G_v\) is \(fv^\epsilon\), and v commutes
with \(\tau\). Hence \(C_{G_v}(\tau)=\langle v\rangle\),
because \(C_F(\tau)=1\). It follows that the entire full cusp
condition is equivalent to one of the literal tests
\[
       d_\lambda=b^\lambda
       \qquad\hbox{or}\qquad
       d_\lambda=v b^\lambda v.
 \tag{A.86.14}
\]
The converse uses the even conjugators 1 and \(vb^{-1}\),
respectively. Returning to e, the corresponding original F-conjugators
in (A.86.10) are \(q^{-1}\) and \(vb^{-1}q^{-1}\).
The tests cannot both hold: otherwise the nontrivial involution v
would belong to the torsion-free centralizer of \(b^\lambda\).
This does not assert that there are only two unnormalized conjugators;
their procyclic centralizer factors have been retained throughout.

\emph{A45 selects one of the two tests.}\par
A.45 supplies the particular unit d in (A.86.1), its dyadic residue
one modulo four, and the exact equivalence between its cusp equation
and full conjugacy. In the actual homology quotient
\[
       G_v\longrightarrow
       (F_{\mathrm{ab}}/2F_{\mathrm{ab}})\rtimes\langle v\rangle,
\]
write X,Y for the two coordinate vectors. The image of e is \((X,v)\).
The image of q is \((t,t)\), fixed by interchange, so the image of b
is also \((X,v)\). This element has square \((X+Y,1)\) and
order four. Thus
\[
 \overline{b^d}=(X,v),\qquad
 \overline{vb^dv}=(Y,v),\qquad
 \overline{x^dv}=(X,v).
 \tag{A.86.15}
\]
The second test of (A.86.14) is impossible. Consequently the actual
full cusp equation for d is exactly \(x^dv=b^d\), or
\(\Omega_\alpha=1\). Equation (A.86.10) gives (A.86.6).
A.39 and A.45 give the precise equivalences in (A.86.5);
the descents inside \(\widehat{GT}_0\) and
\(\widehat{GT}\) are those of A.6 and A.36. Once the full
cusp equation holds, C.29 makes the normalized substitution
\(x\mapsto x^d,\ y\mapsto w^d\) an automorphism.
Unit powering does not change \(\overline{\langle x,w\rangle}\),
so A.85 implies the full Fox-unit conclusion. This is a sufficient
implication from the marked cusp equation; its converse is not used.

\emph{The exact C.44 solution and its all-head unit.}\par
Let \(j=qvq^{-1}\). It has order two and commutes with r because
v commutes with \(\tau\). Both e and j have v-parity one, so
\(g=ej=xu\,v(q)q^{-1}\) lies in F. Equations (A.86.7) and
(A.86.13) give, in the stated order,
\[
       \sigma(g)g^{-1}=(e^{-1}j)(je^{-1})=e^{-2}=z.
 \tag{A.86.16}
\]
The actual involution \(\sigma=\operatorname{Inn}(h_0)\tau\)
has the full regular-pants peripheral action and sends z to its
inverse literally. Thus (A.86.16) is exactly C.44's equation, and
its uniqueness identifies this g with that theorem's actual solution.

We recall the full-unit step to retain its coefficient scope.
Put \(K=\langle\!\langle z\rangle\!\rangle_F\) and
\(A=F/K=\widehat{\mathbb Z}\), with generator the image of x.
C.44's exact compact Alexander calculation gives
\(K/\Phi_2(K)=\mathbb F_2[[A]][z]\). The normalized
\(\sigma\) acts trivially on A and on this module. If
\(\sigma(x)=\kappa x\), \(\theta=[\kappa]\), and
g has image \(a_{\mathrm{reg}}\in A\), equation (A.86.16)
implies
\[
 \theta D_{a_{\mathrm{reg}}}=1,
 \qquad (T-1)D_a=T^a-1
       \quad\hbox{in }\mathbb F_2[[\widehat{\mathbb Z}]].
 \tag{A.86.17}
\]
Multiplication by \(T-1\) identifies this full completed ring
with its augmentation ideal: at the finite cyclic heads its image
is that ideal, and the transition from order \(2m\) to m kills
the constant kernel vectors. Thus the displayed divided differences
are uniquely defined, as in the complete proof of C.44. The first
equation makes \(D_{a_{\mathrm{reg}}}\) a unit. Augmentation
makes \(a_{\mathrm{reg}}\) odd. For any odd prime l, evaluation
of T at a primitive lth root in a finite field of characteristic two
would make that unit zero if \(a_{\mathrm{reg}}=0\pmod l\).
It follows that \(a_{\mathrm{reg}}\) is nonzero modulo every
prime, hence is a full unit. These cyclic heads retain the full
procyclic A; they are not its maximal pro-two quotient.

Finally v interchanges the abelian coordinates and u has zero
abelianization. Hence
\[
 g_{\mathrm{ab}}=(1,0)+(-t,t)-(t,-t)=(1-2t,2t).
\]
Killing z imposes \(y=x^{-1}\), so the actual C.44 unit is
\(a_{\mathrm{reg}}=1-4t\), as claimed. If \(n=(d-1)/2\),
the coefficients of \(b,b^d,x^dv\) in
\(F_{\mathrm{ab}}\rtimes\langle v\rangle\) are respectively
\[
 (1-2t,2t),\qquad(n+1-2t,n+2t),\qquad(d,0).
\]
Their difference proves the last formula in (A.86.4). In particular
\[
       \Omega_\alpha\in F'
       \quad\Longleftrightarrow\quad
       a_{\mathrm{reg}}=d
       \quad\Longleftrightarrow\quad
       t=-\frac{d-1}{4}.
 \tag{A.86.18}
\]
The divisions are defined uniquely in \(\widehat{\mathbb Z}\)
because d is one modulo four at two and its additive group is
torsion-free. No identification is made between the actual dihedral
rotation exponent k and the distinct abelianization coordinate t.

Equation (A.86.18) is only the first coefficient of the full word.
The relative pronilpotent conjugacy in C.44 does not preserve the
three original peripheral markings and does not prove
\(a_{\mathrm{reg}}=d\), much less \(\Omega_\alpha=1\).
Nor is A.42's dyadic square comparison for full
\(\widehat{GT}\)-involutions promoted to an odd-prime or general
\(\widehat{GT}_0\) conclusion. Both original main decisions
remain outside the conclusion of this statement.
\end{proof}
